\chapter{Результаты и их обсуждение}
\label{ch:results}

Для нахождения скорости откачки по формуле (1), нашли объём высоковакуумной части установки (см. Приложение Г). Он составил $V_\text{вв} = 1253.00 \pm 0.11 \ \text{см}^3$ (см. Приложение Б).

Чтобы найти зависимость давления от времени при улучшении вакуума от начального давления до предельного, сначала определили предельное давление $P_\text{пр} = 6.4 
\cdot 10^{-5} \ \text{мм. рт. ст.}$ (см. Приложение В). Затем ухудшили вакуум до давления $P_0 = 5.4 \cdot 10^{-4} \text{мм. рт. ст.}$, запустили откачку высоковакуумного баллона и записали зависимость давления $P$ от времени $t$ в процессе откачки (см. Приложение Б).

Для нахождения скорости откачки по зависимости давления $P$ от времени $t$ в процессе откачки построили график в координатах $ln (P - P_\text{пр})$ и $t$ (см. рис. 1).

% если логарифм от размерной величины, то нужно явно указывать единицы измерения величины.

\begin{figure}[ht]
\center\begin{gnuplot}[terminal = pdf]
    set xl "t, c"
    set yl "ln (P)"
    set xr [0:38]
    set format y "%.1f"
    set grid
    f(x) = k*x + b
    fit f(x) "txt/P(t)improve.txt" u ($1 - 48):(log($2 - 0.64)) via k,b
    plot "txt/P(t)improve.txt" u ($1 - 48):(log($2 - 0.64)):(0.5):(0.05 / ($2)) w xye pt 0 lc "red" lw 2 t "Данные", \ 
    f(x) lw 3 pt 7 ps 0.5 t "Аппроксимация"
\end{gnuplot}  
\vspace{-20pt}
\caption{Зависимость давления от времени при откачке вакуума, давление показано в логарифмическом масштабе. Давление измерялось в мм. рт. ст. %логарифмические координаты, или ln (P), типа плюс какая-то константа.
Красными крестами на графике обозначены данные, полученные из экспериментальных, зелёной линией показана аппроксимация зависимости. Для нахождения скорости откачки аппроксимировали график к линейной зависимости, чтобы описать поведение вакуума при откачке формулой (1).}
\label{fig:lnp(t)}
\end{figure}

Для нахождения скорости откачки аппроксимировали график к линейной зависимости, чтобы описать поведение вакуума при откачке формулой (1). Нашли коэффициент наклона графика, чтобы из него, зная откачиваемый объём, найти скорость откачки $W$. Коэффициент наклона графика получился равным $k = - 0.253 \pm 0.010 \ \text{с}^{-1}$.

Нашли искомую скорость откачки:
\begin{align*}
                W = (3.17 \pm 0.13) \cdot 10^2~\text{см}^3/\text{с}
\end{align*}
Результат показал, что предложенный метод позволяет измерить скорость откачки с погрешностью 4.1 \%.

Также полезно знать величину потока газа $Q_\text{н}$, поступающего из насоса назад в откачиваемую систему при отключенном насосе. Её можно оценить по формуле $P_\text{пред}W = Q_\text{д} + Q_\text{и} + Q_\text{н}$. Величину $(Q_\text{д} + Q_\text{и})$ можно найти, учитывая, что для этого случая $V_\text{вв} dP = (Q_\text{д} + Q_\text{и})dt$ (см. Приложение А). Для определения $Q_\text{н}$ построили график зависимости $P(t)$(см. рис. 2), аппроксимировали его к линейной зависимости и нашли коэффициент наклона полученной прямой $k = (6.4 \pm 1.0) \cdot 10^{-6}~\text{мм. рт. ст.}/ \text{с}$.

\begin{figure}[ht]
\center\begin{gnuplot}[terminal = pdf]
    set xl "t, c"
    set yl "P, мм. рт. ст. · 10^{−4}"
    set format y "%.2f"
    set grid
    set key b
    f(x) = k*x + b
    fit f(x) "txt/P(t)worse.txt" u 1:2 via k,b
    plot "txt/P(t)worse.txt" u 1:2:(0.5):(0.05) w xye pt 0 lc "red" lw 2 t "Данные", \ 
    f(x) lw 3 pt 7 ps 0.5 t "Аппроксимация"
\end{gnuplot}
\vspace{-20pt}
\caption{Полученная экспериментально зависимость давления P от времени t после отключения насоса. График хорошо аппроксимируется к линейной зависимости, это подтверждает справедливость формулы $V_\text{вв} dP = (Q_\text{д} + Q_\text{и})dt$ для описания поведения газа после отключения насоса. Коэффициент наклона графика составил $k = (6.4 \pm 1.0) \cdot 10^{-6}~\text{мм. рт. ст.}/ \text{с}$.}
\label{fig:dT_p}
\end{figure}

Нашли величину $Q_\text{н}$:
\begin{align*}
    Q_\text{н} = (1.23 \pm 0.15) \cdot 10^{-2}~\text{см}^3/\text{с}.
\end{align*}

Чтобы найти производительность диффузионного насоса, ввели в систему со стороны капилляра между высоковакуумной и форвакуумной частями искуственную течь и нашли установившееся при этом давление $P_\text{уст}$ и давление в форвакуумной части $P_\text{фв}$, по параметрам трубки определили её пропускную способность $C_\text{тр}$ (см. Приложение Б). Скорость откачки посчитали по формуле (см. приложение А):

\begin{align*}
            W = \frac{C_\text{тр} P_\text{фв}}{P_\text{уст} - P_\text{пред}}
 \end{align*}
 Получили 
 \begin{align*}
 W = 27 \pm 10 \ \text{см}^3 / \text{с}.
 \end{align*}
 Найденное значение оказалось меньше полученного при помощи анализа зависимости P(t) в процессе откачки. Это можно обосновать тем, что после создания течи эффективная пропускная способность системы вместе с капилляром снизилась, что не позволило достичь максимальной производительности насоса. Из этого следует, что скорость откачки зависит также от пропускной способности трубок, потому что скорость откачки через узкий каппиляр оказалась меньше, чем через более широкую трубку.